\chapter*{Conclusion generale}
\addcontentsline{toc}{chapter}{Conclusion generale}

Ce projet de fin d'etudes avait pour ambition de transformer un besoin metier concret en une solution numerique complete. L'objectif a ete atteint avec la realisation de \textbf{\ProjectOfficialTitle}, ci-apres \textbf{\ProjectName}, une plateforme integree de veille des prix internationaux, allant de la collecte web jusqu'a la visualisation decisionnelle.

\section*{Bilan des apports}
Les principaux apports du projet sont :
\begin{itemize}[leftmargin=1.2cm]
	\item automatisation d'un processus historiquement lourd et partiellement manuel ;
	\item centralisation des donnees prix dans un modele relationnel fiable ;
	\item normalisation des prix en EUR pour des comparaisons internationales ;
	\item disponibilite d'un dashboard professionnel pour les profils metier et admin.
\end{itemize}

\section*{Contribution des etudiants}
Le projet valorise une collaboration reelle entre \textbf{Khalil Amamri} et \textbf{Mahdi ElhadjAmor}, avec une implication continue sur la conception, l'implementation et la validation. Cette dynamique d'equipe a permis d'assurer la coherence de la solution de bout en bout.

\section*{Impact pour l'entreprise}
Pour CHO, la plateforme apporte une meilleure visibilite sur les prix, une reduction du temps d'analyse et une meilleure capacite de pilotage des decisions commerciales dans un environnement international.

\section*{Perspectives}
Les perspectives naturelles d'evolution sont :
\begin{itemize}[leftmargin=1.2cm]
	\item integration de nouveaux pays, enseignes et familles de produits ;
	\item mise en place d'alertes automatiques (variation forte, absence de donnees, anomalies) ;
	\item industrialisation des tests automatiques et deploiement CI/CD.
\end{itemize}

En conclusion, \ProjectName{} constitue une base solide et evolutive pour la veille strategique des prix, avec une valeur immediate pour les utilisateurs et un potentiel important pour les futures iterations.

\clearpage
